\paragraph{bin-fold 群胚代数的忠实表示尺度、PI 阈值与六倍窗半尺度缺陷}
在 bin-fold 纤维群胚代数的 Morita--\(K\) 理论、中心计数、规范群交换化以及六倍窗零余类放大律都已建立之后，还可继续从同一份 Wedderburn 退化直方图抽取出下列五条对象层后果。它们分别刻画最小忠实表示尺度、Morita 类型对块大小的遗忘、PI 阈值的最大纤维判别、六倍窗信息缺口的黄金半尺度锁定，以及 window-\(6\) 上三种彼此不重合的刚性尺度。

\begin{theorem}[bin-fold 纤维群胚代数的最小忠实表示维数等于全部微态数]\label{thm:xi-time-part60aca-minimal-faithful-dimension-microstate-count}
\leanverified{paper\_xi\_time\_part60aca\_minimal\_faithful\_dimension\_microstate\_count}
沿用定义 \ref{def:fold-bin-fiber-groupoid} 与命题 \ref{prop:fold-bin-groupoid-algebra-wedderburn} 的记号，记
\[
A_m:=\CC[\mathcal R_m^{\mathrm{bin}}]
\cong
\bigoplus_{w\in X_m}M_{d_w}(\CC),
\qquad
d_w:=d_m^{\mathrm{bin}}(w).
\]
定义
\[
\mu_{\mathrm{faith}}(A_m)
:=
\min\Bigl\{
\dim_{\CC}H:\ 
\exists\ \rho:A_m\hookrightarrow \operatorname{End}_{\CC}(H)\ \text{为忠实幺代数表示}
\Bigr\}.
\]
则有严格恒等式
\[
\mu_{\mathrm{faith}}(A_m)
=
\sum_{w\in X_m} d_w
=
2^m.
\]
等价地，\(A_m\) 的最小忠实 Hilbert 载体维数恰等于全部微态总数。
\end{theorem}
\begin{proof}
对每个 \(w\in X_m\)，记 \(e_w\in A_m\) 为对应于 \(M_{d_w}(\CC)\) 这一 Wedderburn 块的中心极小幂等元。它们满足
\[
e_we_{w'}=\delta_{w,w'}e_w,
\qquad
\sum_{w\in X_m}e_w=1.
\]
任取忠实幺表示
\[
\rho:A_m\hookrightarrow \operatorname{End}_{\CC}(H).
\]
由于 \(\rho\) 单射，必有 \(\rho(e_w)\neq 0\) 对一切 \(w\) 成立。令
\[
H_w:=\rho(e_w)H.
\]
由幂等元两两正交且和为 \(1\)，得到直和分解
\[
H=\bigoplus_{w\in X_m}H_w.
\]

限制到第 \(w\) 块，得到一个非零表示
\[
M_{d_w}(\CC)\cong e_wA_m\longrightarrow \operatorname{End}_{\CC}(H_w).
\]
由于 \(M_{d_w}(\CC)\) 是简单代数，非零表示的核只能是零，故该表示自动忠实。另一方面，有限维 \(M_{d_w}(\CC)\)-模必分解为标准模 \(\CC^{d_w}\) 的若干重直和，因此
\[
\dim_{\CC}H_w\ge d_w.
\]
求和即得
\[
\dim_{\CC}H=\sum_{w\in X_m}\dim_{\CC}H_w\ge \sum_{w\in X_m}d_w.
\]

反向地，对每个 \(w\) 取标准表示
\[
M_{d_w}(\CC)\curvearrowright \CC^{d_w},
\]
再作外直和
\[
\bigoplus_{w\in X_m}\CC^{d_w},
\]
便得到 \(A_m\) 的忠实表示，其维数恰为 \(\sum_w d_w\)。因此
\[
\mu_{\mathrm{faith}}(A_m)=\sum_{w\in X_m}d_w.
\]
最后，\(\Omega_m=\{0,1\}^m\) 被各条纤维
\[
\Fold_m^{\mathrm{bin}}{}^{-1}(w)
\qquad (w\in X_m)
\]
不交分割，故
\[
\sum_{w\in X_m}d_w=|\Omega_m|=2^m.
\]
证毕。
\end{proof}

\leanverified{paper\_xi\_time\_part60aca\_morita\_type\_forgets\_block\_sizes}
\begin{theorem}[bin-fold 纤维群胚代数的 Morita 类型完全坍缩到 Fibonacci 中心秩]\label{thm:xi-time-part60aca-morita-type-forgets-block-sizes}
沿用定理 \ref{thm:xi-time-part60aca-minimal-faithful-dimension-microstate-count} 的记号。则 \(A_m\) 与
\[
\CC^{F_{m+2}}
\]
Morita 等价。进一步，在自然的 \(\ast\)-结构下将 \(A_m\) 视为有限维 \(C^\ast\)-代数，则其有序 \(K_0\) 群满足
\[
K_0(A_m)\cong \ZZ^{F_{m+2}},
\qquad
K_0(A_m)^+\cong \NN^{F_{m+2}}.
\]
因此，bin-fold 纤维群胚代数的 Morita 类型完全忘记了块大小 \((d_w)\)，只保留块数 \(|X_m|=F_{m+2}\)。
\end{theorem}
\begin{proof}
Morita 等价
\[
A_m\ \text{与}\ \CC^{F_{m+2}}
\]
以及群同构
\[
K_0(A_m)\cong \ZZ^{F_{m+2}}
\]
已由定理 \ref{thm:xi-foldbin-groupoid-morita-k-center-dimension-lb} 给出。现只需说明正锥。

由 Wedderburn 分解，
\[
A_m\cong \bigoplus_{w\in X_m}M_{d_w}(\CC)
\]
作为有限维 \(C^\ast\)-代数成立。于是
\[
K_0(A_m)\cong \bigoplus_{w\in X_m}K_0\bigl(M_{d_w}(\CC)\bigr)
\cong \ZZ^{|X_m|},
\]
其中每个直和因子由相应矩阵块中的投影类生成。对矩阵代数 \(M_d(\CC)\)，投影类恰按其秩由 \(\NN\) 参数化，因此正锥逐块相加后得到
\[
K_0(A_m)^+\cong \NN^{|X_m|}.
\]
再由 \(|X_m|=F_{m+2}\)，遂得
\[
K_0(A_m)^+\cong \NN^{F_{m+2}}.
\]
证毕。
\end{proof}

\begin{theorem}[bin-fold 纤维群胚代数的 PI 阈值由最大纤维完全决定]\label{thm:xi-time-part60aca-pideg-maxfiber-threshold}
\leanverified{paper\_xi\_time\_part60aca\_pideg\_maxfiber\_threshold}
沿用定理 \ref{thm:xi-time-part60aca-minimal-faithful-dimension-microstate-count} 的记号，并定义
\[
D_m^{\mathrm{bin}}
:=
\max_{w\in X_m}d_m^{\mathrm{bin}}(w)
=
d_{\max}^{\mathrm{bin}}(m).
\]
对 \(r\ge 1\)，记 \(2r\) 次标准多项式为
\[
s_{2r}(x_1,\dots,x_{2r})
:=
\sum_{\sigma\in \mathfrak S_{2r}}
\operatorname{sgn}(\sigma)\,
x_{\sigma(1)}\cdots x_{\sigma(2r)}.
\]
则
\[
\operatorname{PIdeg}(A_m)=D_m^{\mathrm{bin}}.
\]
等价地，
\[
A_m\models s_{2D_m^{\mathrm{bin}}},
\qquad
A_m\not\models s_{2D_m^{\mathrm{bin}}-2}.
\]
特别地，当 \(m=6\) 时，由直方图 \(2:8,\ 3:4,\ 4:9\) 得
\[
\operatorname{PIdeg}(A_6)=4,
\qquad
A_6\models s_8,
\qquad
A_6\not\models s_6.
\]
\end{theorem}
\begin{proof}
由 Wedderburn 分解，
\[
A_m\cong \bigoplus_{w\in X_m}M_{d_w}(\CC).
\]
对有限维半单复代数
\[
A\cong \bigoplus_{j=1}^{r}M_{n_j}(\CC),
\]
一个多项式恒等式在 \(A\) 上成立，当且仅当它在每个矩阵块 \(M_{n_j}(\CC)\) 上都成立。由 Amitsur--Levitzki 定理，
\[
M_{n_j}(\CC)\models s_{2n_j},
\qquad
M_{n_j}(\CC)\not\models s_{2n_j-2}.
\]
因此 \(A\) 的 PI 度数恰为 \(\max_j n_j\)。将
\[
n_j=d_w
\]
代入即得
\[
\operatorname{PIdeg}(A_m)=\max_{w\in X_m}d_w=D_m^{\mathrm{bin}}.
\]
等价的标准恒等式表述随即成立。

当 \(m=6\) 时，推论 \ref{cor:fold-bin-m6-fiber-hist} 给出不同块尺寸恰为 \(2,3,4\)，故 \(D_6^{\mathrm{bin}}=4\)。于是
\[
\operatorname{PIdeg}(A_6)=4,
\qquad
A_6\models s_8,
\qquad
A_6\not\models s_6.
\]
证毕。
\end{proof}

\begin{theorem}[六倍窗共振把碰撞与 KL 缺口锁定在黄金半尺度]\label{thm:xi-time-part60aca-sixfold-halfscale-defect}
沿用定理 \ref{thm:xi-fold-zero-arithmetic-subsequence-amplification} 的记号，特别记
\[
p_m(r):=\frac{d_m(r)}{2^m},
\qquad
u_m(r):=\frac{1}{F_{m+2}}.
\]
若
\[
m+2\equiv 0\pmod 6,
\qquad
d:=\frac{m+2}{2},
\]
则 bin-fold 的碰撞缺口与 KL 缺口满足
\[
1-\mathsf{Col}_m
\ge
\frac{F_d}{F_{m+2}}
=
\frac{1}{L_d}
=
\varphi^{-d}\bigl(1+O(\varphi^{-2d})\bigr),
\]
以及
\[
\log F_{m+2}-D_{\mathrm{KL}}(p_m\|u_m)
\ge
-\log\!\left(1-\frac{F_d}{F_{m+2}}\right)
=
\varphi^{-d}\bigl(1+O(\varphi^{-d})\bigr).
\]
等价地，
\[
1-\mathsf{Col}_m
\ge
\varphi^{-(m+2)/2}\bigl(1+O(\varphi^{-m})\bigr),
\]
\[
\log F_{m+2}-D_{\mathrm{KL}}(p_m\|u_m)
\ge
\varphi^{-(m+2)/2}\bigl(1+O(\varphi^{-m/2})\bigr).
\]
因此，在六倍窗子序列上，bin-fold 对均匀分布的接近程度不可能优于一个黄金半尺度 \(\varphi^{-m/2}\)。
\end{theorem}
\begin{proof}
由结论 \ref{con:xi-fold-window6-lucas-barrier-collision-gap}，
\[
1-\mathsf{Col}_m
\ge
\frac{F_d}{F_{2d}}
=
\frac{1}{L_d},
\]
其中 \(L_d\) 为第 \(d\) 个 Lucas 数。由 Binet 公式，
\[
L_d=\varphi^d+(-\varphi)^{-d}
=
\varphi^d\bigl(1+O(\varphi^{-2d})\bigr),
\]
从而
\[
\frac{1}{L_d}
=
\varphi^{-d}\bigl(1+O(\varphi^{-2d})\bigr).
\]
又因 \(2d=m+2\)，故第一条渐近式可改写为
\[
1-\mathsf{Col}_m
\ge
\varphi^{-(m+2)/2}\bigl(1+O(\varphi^{-m})\bigr).
\]

另一方面，条件 \(m+2\equiv 0\pmod 6\) 等价于 \(m\equiv 4\pmod 6\)。于是定理 \ref{thm:xi-fold-zero-arithmetic-subsequence-amplification} 给出
\[
D_{\mathrm{KL}}(p_m\|u_m)
\le
\log\!\bigl(F_{m+2}-F_d\bigr).
\]
因此
\[
\log F_{m+2}-D_{\mathrm{KL}}(p_m\|u_m)
\ge
\log F_{m+2}-\log\!\bigl(F_{m+2}-F_d\bigr)
=
-\log\!\left(1-\frac{F_d}{F_{m+2}}\right).
\]
再由
\[
\frac{F_d}{F_{m+2}}
=
\frac{F_d}{F_{2d}}
=
\frac{1}{L_d}
=
\varphi^{-d}\bigl(1+O(\varphi^{-2d})\bigr),
\]
以及
\[
-\log(1-x)=x+O(x^2)
\qquad(x\to 0),
\]
便得到
\[
\log F_{m+2}-D_{\mathrm{KL}}(p_m\|u_m)
\ge
\varphi^{-d}\bigl(1+O(\varphi^{-d})\bigr).
\]
把 \(d=(m+2)/2\) 代回即得所示第二条渐近式。证毕。
\end{proof}

\begin{theorem}[window-\texorpdfstring{$6$}{6} 上 Morita 秩、PI 阈值与忠实表示维数的三尺度分裂]\label{thm:xi-time-part60aca-window6-triple-rigidity-scales}
取 \(m=6\)。则
\[
A_6
\cong
M_2(\CC)^{\oplus 8}\oplus M_3(\CC)^{\oplus 4}\oplus M_4(\CC)^{\oplus 9}.
\]
从而同一非交换对象上同时出现三条互不重合的刚性尺度：
\[
\#\operatorname{Irr}(A_6)
=
\operatorname{rank}_{\ZZ}K_0(A_6)
=
21,
\qquad
\operatorname{PIdeg}(A_6)=4,
\qquad
\mu_{\mathrm{faith}}(A_6)=64.
\]
此外，
\[
A_6\models s_8,
\qquad
A_6\not\models s_6,
\qquad
\bigl(G_6^{\mathrm{bin}}\bigr)^{\mathrm{ab}}\cong (\ZZ/2\ZZ)^{21}.
\]
因此，window-\(6\) 上的 Morita--\(K\) 理论尺度、PI 尺度与忠实表示尺度并不塌缩为同一数值。
\end{theorem}
\begin{proof}
由推论 \ref{cor:fold-bin-m6-fiber-hist}，
\[
\#\{w\in X_6:\ d_6^{\mathrm{bin}}(w)=2\}=8,
\qquad
\#\{w\in X_6:\ d_6^{\mathrm{bin}}(w)=3\}=4,
\qquad
\#\{w\in X_6:\ d_6^{\mathrm{bin}}(w)=4\}=9.
\]
代入 Wedderburn 分解便得
\[
A_6
\cong
M_2(\CC)^{\oplus 8}\oplus M_3(\CC)^{\oplus 4}\oplus M_4(\CC)^{\oplus 9}.
\]

由定理 \ref{thm:xi-time-part60aca-morita-type-forgets-block-sizes}，
\[
\#\operatorname{Irr}(A_6)
=
\operatorname{rank}_{\ZZ}K_0(A_6)
=
|X_6|
=
21.
\]
由定理 \ref{thm:xi-time-part60aca-pideg-maxfiber-threshold}，
\[
\operatorname{PIdeg}(A_6)=4,
\qquad
A_6\models s_8,
\qquad
A_6\not\models s_6.
\]
再由定理 \ref{thm:xi-time-part60aca-minimal-faithful-dimension-microstate-count}，
\[
\mu_{\mathrm{faith}}(A_6)=2^6=64.
\]
等价地，直接由直方图求和也可写成
\[
\mu_{\mathrm{faith}}(A_6)
=
8\cdot 2+4\cdot 3+9\cdot 4
=
64.
\]
最后，定理 \ref{thm:xi-foldbin-gauge-group-order-abel-center-closed} 给出
\[
\bigl(G_6^{\mathrm{bin}}\bigr)^{\mathrm{ab}}\cong (\ZZ/2\ZZ)^{21},
\]
故其 \(\FF_2\)-维数与第一条尺度一致，而与后两条尺度不同。证毕。
\end{proof}

\paragraph{bin-fold 群胚代数的有序 \texorpdfstring{$K_0$}{K0}、稳定刚性与自同构阈值}
在 bin-fold 纤维群胚代数的 Wedderburn 信息完备性、Morita--\(K\) 理论、PI 阈值与 window-\(6\) 的异常资源分叉已经建立之后，还可把同一退化谱直方图继续压缩为如下六条对象层结论：带单位元类的有序 \(K_0\) 完全分类、稳定同构刚性、PI 度数的黄金下界、自同构群的半直积分解、中心块数与异常比特数的严格对齐，以及首次非交换阈值的精确定位。

\begin{theorem}[bin-fold 纤维群胚代数的有序 \texorpdfstring{$K_0$}{K0} 完全分类]\label{thm:xi-time-part60acb-binfold-ordered-k0-complete-classification}
沿用定义 \ref{def:fold-bin-fiber-groupoid} 与命题 \ref{prop:fold-bin-groupoid-algebra-wedderburn} 的记号，记
\[
A_m:=\CC[\mathcal R_m^{\mathrm{bin}}],
\qquad
\mathcal R_m^{\mathrm{bin}}
=
\bigsqcup_{w\in X_m}(F_w\times F_w),
\qquad
d_w:=|F_w|=d_m^{\mathrm{bin}}(w).
\]
则在由 Wedderburn 块给出的自然坐标下，有
\[
\boxed{\;
\bigl(K_0(A_m),K_0(A_m)^+,[1_{A_m}]\bigr)
\cong
\bigl(\ZZ^{X_m},\NN^{X_m},(d_w)_{w\in X_m}\bigr).\;}
\]
特别地，带 order unit 的有序 \(K_0\) 唯一恢复纤维多重度多重集
\[
\{d_w:w\in X_m\},
\]
反过来，该多重度多重集也唯一决定 \(A_m\) 的复代数同构类。
\end{theorem}
\begin{proof}
由命题 \ref{prop:fold-bin-groupoid-algebra-wedderburn}，
\[
A_m\cong \bigoplus_{w\in X_m}M_{d_w}(\CC).
\]
因此
\[
K_0(A_m)\cong \bigoplus_{w\in X_m}K_0\!\bigl(M_{d_w}(\CC)\bigr).
\]
对每个矩阵块 \(M_{d_w}(\CC)\)，以秩一投影类作为 \(K_0\)-生成元，则
\[
K_0\!\bigl(M_{d_w}(\CC)\bigr)\cong \ZZ,
\qquad
K_0\!\bigl(M_{d_w}(\CC)\bigr)^+\cong \NN,
\qquad
[1_{M_{d_w}(\CC)}]=d_w.
\]
逐块直和便得到
\[
K_0(A_m)\cong \ZZ^{X_m},
\qquad
K_0(A_m)^+\cong \NN^{X_m},
\qquad
[1_{A_m}]=(d_w)_{w\in X_m}.
\]

由于 \(\NN^{X_m}\) 的极端射线恰对应于各个 Wedderburn 块，故在不固定块顺序时，order unit 的坐标仍唯一恢复多重度多重集 \(\{d_w\}_{w\in X_m}\) 仅差一个置换。反向方向由定理 \ref{thm:xi-time-part60aaa-wedderburn-information-completeness} 立即给出：纤维直方图，等价地多重度多重集，唯一决定 \(A_m\) 的半单代数同构型。证毕。
\end{proof}

\begin{corollary}[Morita 压缩与稳定刚性]\label{cor:xi-time-part60acb-binfold-morita-stable-rigidity}
设 \(\mathcal K\) 为可分无限维 Hilbert 空间上的紧算子代数。则对每个 \(m\ge 1\)，有
\[
\boxed{\;
A_m\ \text{与}\ \CC^{F_{m+2}}\ \text{Morita 等价}.\;}
\]
并且对任意 \(m,m'\ge 1\)，成立等价
\[
\boxed{\;
A_m\otimes\mathcal K\cong A_{m'}\otimes\mathcal K
\iff
F_{m+2}=F_{m'+2}
\iff
m=m'.\;}
\]
\end{corollary}
\begin{proof}
第一条正是定理 \ref{thm:xi-time-part60aca-morita-type-forgets-block-sizes} 的内容。

另一方面，由 Wedderburn 分解，
\[
A_m\otimes\mathcal K
\cong
\bigoplus_{w\in X_m}\bigl(M_{d_w}(\CC)\otimes\mathcal K\bigr)
\cong
\bigoplus_{w\in X_m}\mathcal K
\cong
\mathcal K^{\oplus F_{m+2}},
\]
因为 \(M_d(\CC)\otimes\mathcal K\cong \mathcal K\) 对每个 \(d\ge 1\) 都成立。
故若
\[
A_m\otimes\mathcal K\cong A_{m'}\otimes\mathcal K,
\]
则由稳定性
\[
K_0(A_m\otimes\mathcal K)\cong K_0(A_m),
\qquad
K_0(A_{m'}\otimes\mathcal K)\cong K_0(A_{m'}),
\]
再结合定理 \ref{thm:xi-time-part60acb-binfold-ordered-k0-complete-classification}，得到
\[
\ZZ^{F_{m+2}}\cong K_0(A_m)\cong K_0(A_{m'})\cong \ZZ^{F_{m'+2}},
\]
从而
\[
F_{m+2}=F_{m'+2}.
\]
反之，若两者相等，则两边都同构于同一个
\(
\mathcal K^{\oplus F_{m+2}}
\)。

最后，由 Fibonacci 数列严格递增，\(F_{m+2}=F_{m'+2}\) 等价于 \(m=m'\)。证毕。
\end{proof}

\begin{theorem}[PI 度数的纤维极值律与黄金下界]\label{thm:xi-time-part60acb-binfold-pideg-golden-lower-bound}
沿用记号
\[
A_m\cong \bigoplus_{w\in X_m}M_{d_w}(\CC),
\qquad
M_m:=\max_{w\in X_m}d_w,
\qquad
\overline d_m:=\frac{2^m}{F_{m+2}}.
\]
则
\[
\boxed{\;
\operatorname{PIdeg}(A_m)=M_m.\;}
\]
等价地，\(A_m\) 的最小标准多项式恒等式阶数恰为
\[
\boxed{\;
2M_m.\;}
\]
进一步，若 \(C_\varphi>0\) 为定理 \ref{thm:fold-critical-resonance-constant} 的黄金共振常数，则
\[
\boxed{\;
\liminf_{m\to\infty}
\frac{\operatorname{PIdeg}(A_m)}{\overline d_m}
\ge
1+C_\varphi.\;}
\]
\end{theorem}
\begin{proof}
由定理 \ref{thm:xi-time-part60aca-pideg-maxfiber-threshold}，已知
\[
\operatorname{PIdeg}(A_m)=M_m,
\qquad
A_m\models s_{2M_m},
\qquad
A_m\not\models s_{2M_m-2},
\]
故前两条断言成立。

再由定理 \ref{thm:xi-time-part9m-maxfiber-uniformity-gap-chain}，
\[
\liminf_{m\to\infty}\frac{M_m}{\overline d_m}\ge 1+C_\varphi.
\]
把 \(\operatorname{PIdeg}(A_m)=M_m\) 代入即得
\[
\liminf_{m\to\infty}
\frac{\operatorname{PIdeg}(A_m)}{\overline d_m}
\ge
1+C_\varphi.
\]
证毕。
\end{proof}

\begin{theorem}[纤维群胚代数自同构群的半直积分解]\label{thm:xi-time-part60acb-binfold-algebra-automorphism-semidirect}
对每个 \(d\ge 1\) 记
\[
n_d(m):=\#\{w\in X_m:\ d_w=d\}.
\]
则
\[
\boxed{\;
\operatorname{Aut}_{\CC\text{-alg}}(A_m)
\cong
\left(\prod_{\substack{d\ge 1\\ n_d(m)\neq 0}}\operatorname{PGL}_d(\CC)^{\,n_d(m)}\right)
\rtimes
\left(\prod_{\substack{d\ge 1\\ n_d(m)\neq 0}}\mathfrak S_{n_d(m)}\right),\;}
\]
其中右端半直积作用由各 \(\mathfrak S_{n_d(m)}\) 对同尺寸 \(\operatorname{PGL}_d(\CC)\)-因子的置换给出。

特别地，在 \(m=6\) 时，由直方图
\[
2:8,\qquad 3:4,\qquad 4:9,
\]
有
\[
\boxed{\;
\operatorname{Aut}_{\CC\text{-alg}}(A_6)
\cong
\bigl(\operatorname{PGL}_2(\CC)^8\times \operatorname{PGL}_3(\CC)^4\times \operatorname{PGL}_4(\CC)^9\bigr)
\rtimes
\bigl(\mathfrak S_8\times \mathfrak S_4\times \mathfrak S_9\bigr).\;}
\]
并且
\[
\boxed{\;
\pi_0\!\bigl(\operatorname{Aut}_{\CC\text{-alg}}(A_6)\bigr)
\cong
\mathfrak S_8\times \mathfrak S_4\times \mathfrak S_9.\;}
\]
\end{theorem}
\begin{proof}
由定理 \ref{thm:xi-time-part60aaa-wedderburn-information-completeness}，
\[
A_m
\cong
\bigoplus_{\substack{d\ge 1\\ n_d(m)\neq 0}}M_d(\CC)^{\oplus n_d(m)}.
\]
由于 \(M_d(\CC)\not\cong M_e(\CC)\)（\(d\neq e\)）作为简单复代数彼此不同构，任一代数自同构都不能在不同尺寸的块之间混合；故
\[
\operatorname{Aut}_{\CC\text{-alg}}(A_m)
\cong
\prod_{\substack{d\ge 1\\ n_d(m)\neq 0}}
\operatorname{Aut}_{\CC\text{-alg}}\!\bigl(M_d(\CC)^{\oplus n_d(m)}\bigr).
\]

现固定一个 \(d\)，记 \(n:=n_d(m)\) 且
\[
B_d:=M_d(\CC)^{\oplus n}.
\]
任一 \(\CC\)-代数自同构 \(\Phi\) 必将 \(B_d\) 的极小中心幂等元两两置换，从而给出一个排列
\[
\sigma\in \mathfrak S_n.
\]
对每个分量，\(\Phi\) 在 \(M_d(\CC)\) 上诱导一个复代数自同构；由 Skolem--Noether 定理，每个此类自同构都由某个
\[
g_i\in \operatorname{PGL}_d(\CC)
\]
给出。于是
\[
\operatorname{Aut}_{\CC\text{-alg}}(B_d)
\cong
\operatorname{PGL}_d(\CC)^n\rtimes \mathfrak S_n,
\]
其中 \(\mathfrak S_n\) 通过置换 \(n\) 个等尺寸分量作用。
对全部 \(d\) 同时取直积，即得所示总分解。

当 \(m=6\) 时，推论 \ref{cor:fold-bin-m6-fiber-hist} 给出
\[
n_2(6)=8,\qquad n_3(6)=4,\qquad n_4(6)=9,
\]
代入即得显式公式。

最后，\(\operatorname{PGL}_d(\CC)\) 在复拓扑下连通，而 \(\mathfrak S_n\) 离散，故
\[
\pi_0\!\bigl(\operatorname{PGL}_d(\CC)^n\rtimes \mathfrak S_n\bigr)\cong \mathfrak S_n.
\]
对 \(d=2,3,4\) 取直积，便得到
\[
\pi_0\!\bigl(\operatorname{Aut}_{\CC\text{-alg}}(A_6)\bigr)
\cong
\mathfrak S_8\times \mathfrak S_4\times \mathfrak S_9.
\]
证毕。
\end{proof}

\begin{theorem}[window-\texorpdfstring{$6$}{6} 的中心块数与异常比特数严格对齐]\label{thm:xi-time-part60acb-window6-center-anomaly-bit-alignment}
记
\[
A_6:=\CC[\mathcal R_6^{\mathrm{bin}}],
\qquad
G_6^{\mathrm{bin}}:=\Aut(\Fold_6^{\mathrm{bin}}).
\]
则
\[
\boxed{\;
\dim_{\CC}Z(A_6)=21,
\qquad
\dim_{\FF_2}\bigl((G_6^{\mathrm{bin}})^{\mathrm{ab}}\bigr)=21.\;}
\]
更具体地：
\begin{enumerate}
  \item 任何保持全部 simple block 逐块可分的交换块账本，至少需要 \(21\) 个独立复坐标；该下界由
  \[
  Z(A_6)\cong \CC^{21}
  \]
  精确达到。
  \item 在纤维独立置换语义下，任何完备承载全部奇偶异常信息的二元异常账本，至少需要 \(21\) 个独立 \(\FF_2\)-坐标；该下界由
  \[
  (G_6^{\mathrm{bin}})^{\mathrm{ab}}\cong (\ZZ/2\ZZ)^{21}
  \]
  精确达到。
\end{enumerate}
因此，window-\(6\) 上最小块分离交换坐标数与最小异常比特数严格相等，二者都恰为 \(21\)。
\end{theorem}
\begin{proof}
由定理 \ref{thm:xi-foldbin-groupoid-morita-k-center-dimension-lb}，
\[
Z(A_6)\cong \CC^{F_8}=\CC^{21},
\]
故
\[
\dim_{\CC}Z(A_6)=21.
\]
记 \(e_w\in Z(A_6)\) 为与各个 Wedderburn 块对应的 \(21\) 个中心原始幂等元。它们两两正交且非零，并满足
\[
\sum_{w\in X_6}e_w=1.
\]
若一个交换块账本仍要把全部 simple block 逐块分开，则这些 \(e_w\) 在该账本中的像必须仍然两两可区分；特别地，它们必须给出 \(21\) 个两两正交的非零幂等元。任何有限维交换半单复代数中，两两正交非零幂等元的个数不超过其复维数，故此类账本的交换坐标数必至少为 \(21\)。而中心 \(Z(A_6)\) 本身恰以 \(21\) 个坐标达到该下界。

另一方面，由定理 \ref{thm:xi-foldbin-gauge-group-order-abel-center-closed}，
\[
\bigl(G_6^{\mathrm{bin}}\bigr)^{\mathrm{ab}}\cong (\ZZ/2\ZZ)^{21},
\]
故
\[
\dim_{\FF_2}\bigl((G_6^{\mathrm{bin}})^{\mathrm{ab}}\bigr)=21.
\]
再由定理 \ref{thm:xi-window6-oracle-anomaly-coordinate-bifurcation}，在保持纤维独立置换语义不变的前提下，完备异常信息的最小二元承载对象恰需 \(21\) 个独立 \(\FF_2\)-坐标，任何更小的方案都必然破坏该逐纤维直积语义。

故块分离交换坐标数与异常比特数在 window-\(6\) 上严格对齐，并且都以 \(21\) 为最小值。证毕。
\end{proof}

\begin{proposition}[bin-fold 群胚代数的首次非交换阈值]\label{prop:xi-time-part60acb-binfold-first-noncommutative-threshold}
\leanverified{paper\_xi\_time\_part60acb\_binfold\_first\_noncommutative\_threshold}
对每个 \(m\ge 1\)，记
\[
A_m\cong \bigoplus_{w\in X_m}M_{d_w}(\CC),
\qquad
|X_m|=F_{m+2}.
\]
则：
\begin{enumerate}
  \item 对所有 \(m\ge 1\)，\(A_m\) 都不是单代数；
  \item \(A_m\) 可交换当且仅当 \(m=1\)；
  \item 因而 bin-fold 群胚代数的首次非交换阈值恰发生在
  \[
  \boxed{\;m=2.\;}
  \]
\end{enumerate}
\end{proposition}
\begin{proof}
由 \(F_{m+2}\ge 2\)（\(m\ge 1\)）可知
\[
|X_m|=F_{m+2}\ge 2.
\]
因此 \(A_m\) 的 Wedderburn 分解中至少含有两个简单块，从而 \(A_m\) 对所有 \(m\ge 1\) 都不是单代数。

再看交换性。有限维半单复代数
\[
\bigoplus_{w\in X_m}M_{d_w}(\CC)
\]
可交换，当且仅当每个块都为 \(1\times 1\)，亦即
\[
d_w=1
\qquad (\forall\,w\in X_m).
\]
这等价于
\[
\sum_{w\in X_m}d_w=|X_m|.
\]
而左端恒等于 \(2^m\)，右端恒等于 \(F_{m+2}\)，故交换性等价于
\[
2^m=F_{m+2}.
\]
当 \(m=1\) 时，
\[
2^1=2=F_3,
\]
于是 \(A_1\cong \CC\oplus \CC\) 确为交换。
当 \(m\ge 2\) 时，Fibonacci 数满足
\[
F_{m+2}<2^m,
\]
例如可由归纳得到：\(F_4=3<4=2^2\)，且若
\[
F_{n+2}<2^n,\qquad F_{n+1}<2^{n-1},
\]
则
\[
F_{n+3}=F_{n+2}+F_{n+1}<2^n+2^{n-1}<2^{n+1}.
\]
从而平均块大小
\[
\frac{1}{|X_m|}\sum_{w\in X_m}d_w=\frac{2^m}{F_{m+2}}>1.
\]
故必存在某个 \(w\) 使 \(d_w\ge 2\)，于是 \(A_m\) 非交换。

综上，\(A_m\) 可交换当且仅当 \(m=1\)，故首次非交换阈值恰为 \(m=2\)。证毕。
\end{proof}

\endinput


\noindent
由此，bin-fold 纤维群胚代数上的对象层刚性已进一步闭合：Morita--\(K\) 理论只保留 Fibonacci 块数，而带单位元类的有序 \(K_0\) 又把全部块大小逐坐标恢复；稳定范畴中幸存的不变量仍只有块数，并因此对分辨率给出严格刚性；PI 理论则精确读取最大纤维，并继承黄金共振常数所强制的统一下界；同时，代数自同构群按等尺寸矩阵块分裂为内自同构 Lie 群与离散置换群的半直积，window-\(6\) 上的中心块数与异常比特数又同时锁定为 \(21\)；最后，非交换性并非渐近出现，而是在 \(m=2\) 处发生首次精确阈值跳变。因而，退化谱直方图不再只是静态组合摘要，而是同时支配有序 \(K\) 理论、稳定类型、PI 阈值、自同构群与非交换起始尺度的统一刚性对象。

\endinput
